% Lösungen Einheit 3 von 4 – Normalform (zusammenbau v0.1)
\einheitenkopf{Einheit 3 von 4 · Normalform}

%% TODO Lösung der Erklärzeile 15g) fehlt in der Bank
\erg{15}{a) Normalform \quad b) $f(0) = 5$; Schnittpunkt $(0|5)$ \quad c) $f(0) = 7$; Schnittpunkt $(0|7)$ \quad d) $f(0) = 4$; Schnittpunkt $(0|4)$ \quad e) $f(0) = 9$; Schnittpunkt $(0|9)$ \quad f) $f(0) = 8$; Schnittpunkt $(0|8)$ \quad g) \ldots \quad h) $f(-3) = 9 - 12 - 2 = -5$ \quad i) $f(x) = x^2 - 4x + 9$, denn $(x - 2)^2 + 5 = x^2 - 4x + 4 + 5$ \quad j) $f(x) = x^2 + 8x + 17$, denn $(x + 4)^2 + 1 = x^2 + 8x + 16 + 1$ \quad k) $(x + 5)^2 - 8 = x^2 + 10x + 25 - 8 = x^2 + 10x + 17$ (wA) – die Behauptung stimmt. \quad l) $S(2|{-3})$ \quad m) $S(3|{-4})$; $f(x) = (x - 3)^2 - 4$; Probe: $(0 - 3)^2 - 4 = 5$ (wA) \quad n) wahr – $f(2) = 4 + 4 - 3 = 5$; wahr – $S(-1|{-4})$; falsch – der Schnittpunkt ist $(0|{-3})$ \quad o) $S(4|{-3})$; $f(x) = (x - 4)^2 - 3$ \quad p) $(x + 4)^2 - 3 = x^2 + 8x + 16 - 3 = x^2 + 8x + 13$ (wA)}
\erg{16}{a) Das Mittelglied fehlt, und $(-4)^2$ ist positiv: $(x - 4)^2 + 1 = x^2 - 8x + 16 + 1 = x^2 - 8x + 17$ \quad b) Auf der $y$-Achse ist $x = 0$; dann fallen $x^2$ und $-7x$ weg, übrig bleibt $4$.}
